\documentclass[10pt,a4paper]{article}

% ----------------- Paquetes -------------------%
\usepackage[T1]{fontenc}
\usepackage[utf8]{inputenc}
\usepackage[a4paper,margin=2.5cm]{geometry}
\usepackage{mathptmx}             
\usepackage{changepage} 
\usepackage{setspace}
\usepackage[spanish]{babel}
\usepackage[labelfont=bf,labelsep=space]{caption}
\usepackage{amssymb}
\usepackage{amsmath}
\usepackage{ebgaramond}
\usepackage{placeins}
\usepackage{etoolbox}
\usepackage{setspace}
\usepackage{parskip} 
\usepackage{tcolorbox}
\usepackage{needspace}
\usepackage{xparse}
\usepackage{ocgx2}
\usepackage{hyperref}
\usepackage{hypcap}


%%%%%%%%%%%%%%%%%%%%%%%%%%%%%%%%%%%%%%%%%%%%%%%%%%%%%%%%%%%%%%%%
% --- Fija primero tus valores globales "buenos" ---
\setstretch{1.2}                 % Interlineado global
\setlength{\parskip}{0.7\baselineskip}
\setlength{\parindent}{0pt}
\newlength{\GlobalParskip}
\setlength{\GlobalParskip}{\parskip}

\newcounter{eqnum}[subsection]
\renewcommand{\theeqnum}{\arabic{eqnum}}
\newcounter{alphaitem}
\newcounter{numitem}

%% ------------------- Recuadros----------------------%%
\newcommand{\puntosclave}[1]{%
  \begin{tcolorbox}[
    colframe=black,
    title={Puntos clave},
    before upper={\setlength{\parindent}{0.5cm}\setlength{\parskip}{0.5\baselineskip}}
  ]
  #1
  \end{tcolorbox}
}

\newcommand{\modificaciones}[1]{
  \begin{tcolorbox}[
    colback=white,
    colframe=red,
    title={Modificaciones a realizar},
    before upper={\setlength{\parindent}{0.5cm}\setlength{\parskip}{0.5\baselineskip}}
  ]
  #1
  \end{tcolorbox}
}

%% ------------------- Preguntas TEST----------------------%%
% Opciones: radio (single-choice)
\NewDocumentCommand{\OptRadio}{m m m}{%
  \noindent\CheckBox[radio,name=#1,value=#2,width=1.2ex,height=1.2ex]{}%
  \;\textbf{#2.}~#3\par
}

% Opciones: checkbox (multi-choice)
\NewDocumentCommand{\OptCheck}{m m}{%
  \noindent\CheckBox[name=#1,width=1.2ex,height=1.2ex,bordercolor={0 0 0}]{}%
  \;#2\par
}

% Pregunta single-choice (radio)
% Args: 1=id, 2=enunciado, 3=opciones, 4=solución+justificación, 5=imagen (o {}), 6=origen
\NewDocumentCommand{\PreguntaTestSingle}{m m m m m m}{%
  \par\medskip
  \noindent\textbf{#1.}~#2\par\smallskip

    % Imagen (si no es {})
    \IfBlankTF{#5}{}{%
      \begin{center}
        \includegraphics[width=0.75\linewidth]{#5}
      \end{center}
    }

    \begin{Form}
      #3
    \end{Form}

    \par\smallskip
    \switchocg{sol-#1}{\fbox{Mostrar / ocultar solución}}%
    \hfill{\footnotesize #6}\par\smallskip

    \begin{ocg}{Solución}{sol-#1}{0}
      \noindent #4
    \end{ocg}
  \par\medskip
}

% Pregunta multi-choice (checkboxes)
\NewDocumentCommand{\PreguntaTestMulti}{m m m m m m}{%
  \par\medskip
  \noindent\textbf{#1.}~#2\par\smallskip

    % Imagen (si no es {})
    \IfBlankTF{#5}{}{%
      \begin{center}
        \includegraphics[width=0.75\linewidth]{#5}
      \end{center}
    }
    \begin{Form}
      #3
    \end{Form}

    \par\smallskip
    \switchocg{sol-#1}{\fbox{Mostrar / ocultar solución}}%
    \hfill{\footnotesize #6}\par\smallskip

    \begin{ocg}{Solución}{sol-#1}{0}
      \noindent #4
    \end{ocg}
  \par\medskip
}

%% ------------------- Listas----------------------%%
    % --- Lista alfabética ---
    \newenvironment{la}
      {\begin{list}{\alph{alphaitem})}{%
          \usecounter{alphaitem}%
          \setlength{\leftmargin}{1cm}%
          \setlength{\parskip}{3pt}% 
          \setlength{\itemsep}{0pt}% ← espacio entre ítems
          \setlength{\parsep}{0pt}%  ← párrafos dentro de un ítem
      }}
      {\end{list}}
      
    % --- Lista numérica ---
    \newenvironment{lnum}
      {\begin{list}{\arabic{numitem}.}{%
          \usecounter{numitem}%
          \setlength{\leftmargin}{1cm}%
          \setlength{\parskip}{3pt}% 
          \setlength{\itemsep}{0pt}% ← espacio entre ítems
          \setlength{\parsep}{0pt}%  ← párrafos dentro de un ítem
      }}
      {\end{list}}
      
% ------------------- Imágenes----------------------%
    \graphicspath{{Img/}}% Rutas por defecto 
    % --- Imagen adaptativa ---
    % Registros de dimensión definidos una sola vez
    \newdimen\imgcap
    \newdimen\imgmin
    \newdimen\imgmax
    \newdimen\imgremain
    \newdimen\imgh
    \newdimen\imgneed
    
    \newcommand{\imagenfit}[3]{%
      \addvspace{10pt}% separación antes
      \par\begingroup
        % Parámetros de composición (asignaciones, no \newdimen)
        \imgcap=1\baselineskip            % alto estimado de título+fuente
        \imgmin=0.15\textheight
        \imgmax=0.15\textheight
        \imgremain=\pagegoal
        \ifdim\pagegoal=\maxdimen \imgremain=0pt\fi
        \advance\imgremain by -\pagetotal
        \advance\imgremain by -\imgcap
        % Decisión de altura destino
        \ifdim\imgremain<\imgmin
          \newpage
          \imgh=0.20\textheight
        \else
          \imgh=\imgremain
          \ifdim\imgh>\imgmax \imgh=\imgmax \fi
        \fi
        % Reservar espacio para evitar cortes del bloque completo
        \imgneed=\imgh \advance\imgneed by \imgcap
        \Needspace{\imgneed}
        % Cuerpo
        \noindent\begin{minipage}{\textwidth}
          \centering
          \captionof{figure}{\textemdash\ #2}\par\vspace{0.3em}% título arriba
          \includegraphics[width=\linewidth,height=\imgh,keepaspectratio]{\detokenize{#1}}\par
          \vspace{0.3em}\small\textit{Fuente: }#3% fuente abajo
        \end{minipage}%
      \endgroup
      \par\addvspace{10pt}%
    }

%------------ Bloque de RECUADRO ESPECIAL -------------%
    \newenvironment{specialblock}[1]{%
      \begin{quote} % crea sangría a izquierda y derecha
      \small % texto más pequeño
      \noindent\textbf{#1}\\[-0.3em] % título en negrita
      \rule{\linewidth}{0.4pt}\\[0.6em]}{
      \end{quote}}
      
%-------------------------- Portada -----------------------%
\newcommand{\oldtitle}[1]{\def\OldTitle{#1}}
\newcommand{\changerationale}[1]{\def\ChangeWhy{#1}}

\newcommand{\changebox}{%
\begin{tcolorbox}[title={Historial de cambios del tema}]
\textbf{Título anterior:} \OldTitle\par
\textbf{Motivo del cambio:} \ChangeWhy
\end{tcolorbox}
}

%-------------------------- Author Font -----------------------%
    \newcommand{\authorfont}[1]{{\fontfamily{EBGaramond-TLF}\selectfont #1}}

%-------------------------- Anexos -----------------------%
    \makeatletter
    \newcounter{anexo}
    \renewcommand{\theanexo}{\Roman{anexo}}
    \newcommand{\anexo}[1]{%
      \clearpage
      \phantomsection
      \refstepcounter{anexo}%
      \noindent\textbf{Anexo \theanexo.- }#1\par\medskip
      \addcontentsline{stc}{section}{Anexo \theanexo.- #1}%
    }
    \makeatother

    %-------------------------- Lista de figuras (personal) -----------------------%
\makeatletter
\newcommand{\MyListOfFigures}{%
  \clearpage
  \phantomsection
  {\centering\bfseries\Large Índice de figuras\par}%
  \medskip
  % Entrada en nuestro índice de contenido (stc)
  \addcontentsline{stc}{section}{Índice de figuras}%
  % Recuperación del listado (sin cabecera propia)
  \begingroup
    \parindent=0pt\parskip=2pt
    \@starttoc{lof}%
  \endgroup
}
\makeatother

%-------------------------- Bibliografía -----------------------%
\makeatletter
% Contador para las referencias
\newcounter{myref}
% Cabecera de la bibliografía, entra en el índice 'stc'
\newcommand{\MyBibliography}{%
  \clearpage
  \phantomsection
  {\centering\bfseries\Large Bibliografía y referencias\par}%
  \medskip
  \addcontentsline{stc}{section}{Bibliografía y referencias}%
  \setcounter{myref}{0}%
}
\newcommand{\MyBibItem}[4]{%
  \refstepcounter{myref}%
  \noindent[\themyref]\ #1.\ \emph{#2}. #3. #4\par
}
\makeatother

%--------- Establecer cuatro niveles de índice --------%
    \setcounter{tocdepth}{4}   
    \setcounter{secnumdepth}{4}% numera hasta \paragraph

%%%%%%%%%%%%%%%%%%%%%%%%%%%%%%%%

\makeatletter

    %--------- Interlineado Global --------%
    \edef\GlobalBaseLineStretch{\baselinestretch}
    
    %--------- Espaciado de seccinones --------%
    \renewcommand\section{\@startsection{section}
      {1}{\z@}
      {2.5ex \@plus 1ex \@minus .2ex} % espacio antes
      {1.5ex \@plus .2ex}             % espacio después
      {\normalfont\Large\bfseries}    % formato del título
    }
    \renewcommand\subsection{\@startsection{subsection}{2}{\z@}
      {3ex \@plus 0.8ex \@minus .2ex}
      {1.5ex \@plus .2ex}
      {\normalfont\large\bfseries}
    }
    \renewcommand\subsubsection{\@startsection{subsubsection}{3}{\z@}
      {3ex \@plus 0.5ex \@minus .2ex}
      {1ex \@plus .2ex}
      {\normalfont\normalsize\bfseries}
    }
    \renewcommand\paragraph{\@startsection{paragraph}{4}{\z@}%
    {-2ex \@plus -0.5ex \@minus -.2ex}% espacio antes
    {0.8ex \@plus .2ex}% espacio después
    {\normalfont\normalsize\itshape}}% ← cursiva en lugar de negrita

    %--------- Ecuaciones --------%   
    % Variables globales para almacenar títulos y notas
    \gdef\leq@current@section{}%
    \gdef\leq@current@subsection{}%
    \gdef\leq@last@printed@section{}%
    \gdef\leq@last@printed@subsection{}%
    
    % Comando para añadir nota (invisible en el cuerpo)
    \newcommand{\eqnote}[1]{%
      \addcontentsline{leq}{leqnote}{#1}%
    }
    
    % Comando para ecuaciones
    \newcommand{\eqblock}[2]{%
      % Verificar si necesitamos imprimir el título de sección
      \ifx\leq@current@section\leq@last@printed@section\else
        \addcontentsline{leq}{leqsection}{\leq@current@section}%
        \global\let\leq@last@printed@section\leq@current@section
        \gdef\leq@last@printed@subsection{}% Reset subsección al cambiar sección
      \fi
      % Verificar si necesitamos imprimir el título de subsección
      \ifx\leq@current@subsection\leq@last@printed@subsection\else
        \ifx\leq@current@subsection\@empty\else
          \addcontentsline{leq}{leqsubsection}{\leq@current@subsection}%
          \global\let\leq@last@printed@subsection\leq@current@subsection
        \fi
      \fi
      % Ecuación
      \refstepcounter{eqnum}%
      \begin{equation*}
        #1\tag{E.\theeqnum}
      \end{equation*}%
      \addcontentsline{leq}{leqt}{[E.\theeqnum] -- #2}%
      \addcontentsline{leq}{leqm}{#1}%
    }
    
    %%% ----- Índice de ecuaciones ----------%%%
    % Tamaño para []
    \newcommand{\leqIndexFont}{\normalsize}
    \newcommand{\leqTitleFont}{\tiny}
    \newcommand{\leqNoteFont}{\tiny}
    % Cabecera de sección 
    \newcommand*\l@leqsection[2]{%
      \addvspace{25pt}% Mayor separación antes de nueva sección
      \noindent\leqIndexFont
      \makebox[\linewidth]{\rule{.38\linewidth}{0.4pt}\ \ \textbf{  \MakeUppercase{#1}}\ \ \rule{.38\linewidth}{0.4pt}}%
      \par\vspace{-10pt}
    }
    % Cabecera de subsección 
    \newcommand*\l@leqsubsection[2]{%
      \par\addvspace{15pt}%
      \noindent\leqIndexFont #1
      \par\vspace{-7pt}
    }
    % Título de ecuación 
    \newcommand*\l@leqt[2]{%
    \par\addvspace{10pt}%
      \par\noindent\leqTitleFont #1\par
    }
    % Miniatura de ecuación
    \newcommand*\l@leqm[2]{%
      \vspace{0.3\baselineskip}%
      \par\scriptsize\noindent\hspace*{1.5em}\(#1\)\par%
    }
    % Nota de ecuación 
    \newcommand*\l@leqnote[2]{%
      \vspace{0.2\baselineskip}%
      \par\noindent\leqNoteFont\hspace*{1.5em}\textbf{Nota.--} #1\par\vspace{0.5\baselineskip}%
    }
    % Generación del índice
    \newcommand{\mylistofequations}{%
      \clearpage
      \phantomsection
      % Título visual de la lista (ajústalo a tu gusto):
      {\centering\bfseries\Large Índice de ecuaciones\par}
      % Entrada en nuestro índice 'stc':
      \addcontentsline{stc}{section}{Índice de ecuaciones}%
      % Llamada a tu lista real (usa el nombre exacto de tu macro):
      \@starttoc{leq}%
    }
    
    %--------- Captura de títulos de sección --------%
    \let\leq@original@section\section
    \let\leq@original@subsection\subsection
    % Flag para saber si estamos en una sección asterisco
    \newif\ifLeqStarSection
    \LeqStarSectionfalse
    
    % Redefinir \section CON CONTROL DE ASTERISCO
    \renewcommand{\section}{%
      \@ifstar{\leq@section@star}{\@dblarg\leq@section@nostar}%
    }
    \def\leq@section@star#1{%
      \global\LeqStarSectiontrue
      \gdef\leq@current@section{#1}%
      \gdef\leq@current@subsection{}%
      \leq@original@section*{#1}%
      \global\LeqStarSectionfalse
    }
    \def\leq@section@nostar[#1]#2{%
      \global\LeqStarSectionfalse
      \gdef\leq@current@section{#2}%
      \gdef\leq@current@subsection{}%
      \leq@original@section[#1]{#2}%
    }
    % Redefinir \subsection
    \renewcommand{\subsection}{%
      \@ifstar{\leq@subsection@star}{\@dblarg\leq@subsection@nostar}%
    }
    \def\leq@subsection@star#1{%
      \global\LeqStarSectiontrue
      \gdef\leq@current@subsection{#1}%
      \leq@original@subsection*{#1}%
      \global\LeqStarSectionfalse
    }
    \def\leq@subsection@nostar[#1]#2{%
      \global\LeqStarSectionfalse
      \gdef\leq@current@subsection{#2}%
      \leq@original@subsection[#1]{#2}%
    }

    % ----- Restauración de valores Globales de estilo ----------%
    % Flags
    \newif\ifInFootnote
    \InFootnotefalse
    \pretocmd{\@footnotetext}{\InFootnotetrue}{}{}
    \apptocmd{\@footnotetext}{\InFootnotefalse}{}{}
    % Profundidad de listas (soporta anidadas)
    \newcount\MyListDepth \MyListDepth=0
    \AtBeginEnvironment{la}{\advance\MyListDepth by 1}
    \AfterEndEnvironment{la}{\advance\MyListDepth by -1}
    \AtBeginEnvironment{lnum}{\advance\MyListDepth by 1}
    \AfterEndEnvironment{lnum}{\advance\MyListDepth by -1}
    \AtBeginEnvironment{itemize}{\advance\MyListDepth by 1}
    \AfterEndEnvironment{itemize}{\advance\MyListDepth by -1}
    \AtBeginEnvironment{enumerate}{\advance\MyListDepth by 1}
    \AfterEndEnvironment{enumerate}{\advance\MyListDepth by -1}
    % Restauración diferida: hazlo al INICIO del próximo párrafo real
    \newif\if@pendingRestore
    \@pendingRestorefalse
    \newcommand{\RequestRestore}{\global\@pendingRestoretrue}
    \everypar{%
      \if@pendingRestore
        \global\@pendingRestorefalse
        \ifInFootnote\else
          \ifnum\MyListDepth=0
            \setlength{\parskip}{\GlobalParskip}%
            \setstretch{\GlobalBaseLineStretch}%
          \fi
        \fi
      \fi
    }
    % Al final de bloques:
    \AfterEndEnvironment{equation}{\RequestRestore}
    \AfterEndEnvironment{equation*}{\RequestRestore}
    \AfterEndEnvironment{la}{\RequestRestore}
    \AfterEndEnvironment{lnum}{\RequestRestore}
    % Al final de macros:
    \apptocmd{\eqblock}{\RequestRestore}{}{}
    \apptocmd{\imagenfit}{\RequestRestore}{}{}
    
\makeatother

% ===================== ToC propio con 4 niveles (único bloque) =====================
\makeatletter

% --- Escritores: añadimos una línea al .stc en cada cabecera numerada ---
% Guardamos las versiones actuales para llamar al comportamiento normal
\let\MyTOC@orig@section\section
\let\MyTOC@orig@subsection\subsection
\let\MyTOC@orig@subsubsection\subsubsection
\let\MyTOC@orig@paragraph\paragraph

% SECTION
\renewcommand{\section}{%
  \@ifstar{\MyTOC@section@star}{\@dblarg\MyTOC@section@nostar}%
}
\def\MyTOC@section@star#1{%
  \MyTOC@orig@section*{#1}% sin entrada al .stc
}
\def\MyTOC@section@nostar[#1]#2{%
  \MyTOC@orig@section[{#1}]{#2}% comportamiento normal (escribe al .toc)
  % Escribimos también nuestra línea al .stc (texto con \numberline)
  \addcontentsline{stc}{section}{\protect\numberline{\thesection}#2}%
}

% SUBSECTION
\renewcommand{\subsection}{%
  \@ifstar{\MyTOC@subsection@star}{\@dblarg\MyTOC@subsection@nostar}%
}
\def\MyTOC@subsection@star#1{%
  \MyTOC@orig@subsection*{#1}%
}
\def\MyTOC@subsection@nostar[#1]#2{%
  \MyTOC@orig@subsection[{#1}]{#2}%
  \addcontentsline{stc}{subsection}{\protect\numberline{\thesubsection}#2}%
}

% SUBSUBSECTION
\renewcommand{\subsubsection}{%
  \@ifstar{\MyTOC@subsubsection@star}{\@dblarg\MyTOC@subsubsection@nostar}%
}
\def\MyTOC@subsubsection@star#1{%
  \MyTOC@orig@subsubsection*{#1}%
}
\def\MyTOC@subsubsection@nostar[#1]#2{%
  \MyTOC@orig@subsubsection[{#1}]{#2}%
  \addcontentsline{stc}{subsubsection}{\protect\numberline{\thesubsubsection}#2}%
}

% PARAGRAPH
\renewcommand{\paragraph}{%
  \@ifstar{\MyTOC@paragraph@star}{\@dblarg\MyTOC@paragraph@nostar}%
}
\def\MyTOC@paragraph@star#1{%
  \MyTOC@orig@paragraph*{#1}%
}
\def\MyTOC@paragraph@nostar[#1]#2{%
  \MyTOC@orig@paragraph[{#1}]{#2}%
  % Si no numeras los \paragraph, cambia \theparagraph por texto plano sin \numberline
  \addcontentsline{stc}{paragraph}{\protect\numberline{\theparagraph}#2}%
}

% --- Impresor del índice propio (lee el .stc) ---
\newcommand{\MyTOC}{%
  \par\bigskip
  {\centering\bfseries\Large Índice de contenido\par}%
  \medskip
  \begingroup
    \parindent=0pt\parskip=2pt
    \@starttoc{stc}%
  \endgroup
}

\makeatother
% ===================== fin ToC propio =====================


%%%%%%%%%%%%%%


% --- Datos del tema ---
\title{La gobernanza económica de la Unión Europea y de la zona euro\\
\large El Semestre Europeo. Las reglas fiscales y el Procedimiento de Desequilibrios Macroeconómicos. El Mecanismo Europeo de Estabilidad. La respuesta común de política fiscal a partir de la Covid-19}
\author{Marcelo Ortega}
\date{Julio 2026}
\newcommand{\TemaCode}{3B45}

\oldtitle{
    B.46. El Sistema Europeo de Bancos Centrales y la política monetaria de la Eurozona: objetivos e instrumentos. La gobernanza económica del euro. La Unión Bancaria y otros desarrollos
}
\changerationale{
    Este tema versa sobre la gobernanza económica y la coordinación de la política fiscal en la UE.
}

\begin{document}


\maketitle
\changebox

\newpage

% --- Página de resumen y modificaciones ---
\puntosclave{ %% Escribir puntos clave Apuntes CECO
     
    }
\vspace{1cm}

\modificaciones{
    \textbf{NOTA (04/10/2026):} Revisar: https://fedea.net/evaluacion-de-las-respuestas-fiscales-europeas-al-bloqueo-del-estrecho-de-ormuz/

    \textbf{NOTA (04/10/2026):} Ojo! Algunos temas de economica internacional (realidad y acontecimientos), enlazar con el riesgo que se esta tratando, por ejemplo riesgo financiero, riesgo médico (covid) etc
}

\newpage
\MyTOC
\newpage

\mylistofequations
\clearpage

\MyListOfFigures
\clearpage


\pagenumbering{arabic}
\setcounter{page}{1}


\section{Introducción}



\subsection{Contextualización}


\subsubsection{Relevancia}




\subsubsection{Problemática}



\subsection{Estructura de la exposición}





\clearpage
\section{Coordinación de la Unión: Instrumentos de gobernanza y supervisión}
\puntosclave{
    Necesidad de un marco de gobernanza económica en la UE y especialmente en la Eurozona. 
    
    El Semestre europeo es el ciclo anual de coordinación europea de políticas fiscales y económicas, siendo esta coordinación más estrecha entre países de la Eurozona. 
    
    El Pacto de Estabilidad y Crecimiento (PEC) es el pilar fundamental de la coordinación a nivel fiscal, compuesto por un conjunto de normas concebidas para garantizar que los países de la UE mantengan unas finanzas públicas saneadas y coordinen sus políticas fiscales. 
    
    Por otro lado, el Procedimiento de Desequilibrios Macroeconómicos (PDM) es el fundamento de la coordinación de políticas macroeconórnicas, tiene por objeto identificar, prevenir y abordar los desequilibrios macroeconómicos potencialmente perjudiciales para la estabilidad económica de un país concreto de la UE, de la zona del euro o de la UE en su conjunto. 
    
    El marco de gobernanza económica se encuentra actualmente en revisión y con ello el PEC y el PDM, con el objetivo de simplificar dicho proceso.
}

Los países de la Unión Europea, y en especial los de la zona del euro, coordinan sus políticas económicas y presupuestarias a lo largo del año para garantizar el buen funcionamiento del proyecto europeo. Actuación cuya necesidad se hizo patente ya desde los inicios del proyecto en común, y que dejó de ser optativa tras la introducción del euro, siendo ésta imprescindible para el correcto funcionamiento de un área UEM. 

\textbf{EXPLICAR POR QUÉ: CASO GRECIAS}

La necesidad de una gobernanza económica para los países del euro se halla en el hecho de que el Tratado de Funcionamiento de la Unión Europea establece una clara división de responsabilidades entre los policy-makers nacionales y europeos en el contexto de la Unión Económica y Monetaria. La política monetaria en una unión monetaria como la eurozona es inherentemente indivisible, conducida a nivel supranacional. Por su parte, otras políticas económicas, como la fiscal o las relacionadas con reformas estructurales, siguen siendo competencia de las autoridades nacionales y, por tanto, reflejan sus prioridades políticas individuales. Sin embargo, una política monetaria orientada a la estabilidad de precios no es suficiente para el correcto funcionamiento, sino que, además, son necesarias políticas fiscales sostenibles, así como otra serie de políticas económicas que promuevan la estabilidad financiera, el crecimiento económico y la cohesión social entre los países que la conforman.


\subsection{¿Cómo ha evolucionado este marco de actuación?}
[...]

La crisis financiera y económica que se inició en 2008 evidenció que la Unión necesitaba un modelo de gobernanza más eficaz que la mera coordinación vigente hasta aquel momento. Los cambios, que continúan produciéndose, en el ámbito de la gobernanza económica incluyen una coordinación y vigilancia reforzadas de las políticas presupuestarias y macroeconómicas, así como la instauración de un marco para la gestión de las crisis financieras. De esta forma, el actual marco de gobernanza económica de la Unión Europea tiene por objeto vigilar, prevenir y corregir las tendencias económicas problemáticas que puedan debilitar las economías nacionales o afectar negativamente a otros países de la UE.Este marco de gobernanza económica de la Unión Europea se vertebra sobre el Semestre Europeo,el ciclo de coordinación de las políticas económicas, presupuestarias, sociales y de empleo dentrode la UE. Aunque inicialmente el Semestre Europeo era principalmente un ejercicio económico, ha evolucionado y su proceso integra ahora otros ámbitos políticos. Se centra en los primeros seis meses de cada año, de ahí que se denomine "Semestre". Durante el Semestre Europeo, los Estados miembros ajustan sus políticas presupuestarias y económicas a las normas acordadas a escala de la UE.

Los procedimientos de coordinación de las políticas económicas y sociales que existían antes de 2010 se aplicaban de forma independiente entre ellos.\footnote{En qué consistían: (i) Cooperación mediante intercambio de información; (ii) Coordinación como gestión de crisis, por ejemplo, la instauración del MEEF en mayo de 2010; y, (iii) delegación de una política a unainstitución europea (Política Monetaria y Política de la Competencia); etc.

Un ejemplo de la escasa voluntad de coordinación fue la que se dió en los primeros años de la crisis sobre las medidas para la recapitalización bancaria. Se abordó el problema desde un enfoque puramente nacional, se puso en entredicho el mercado único en materia financiera y el papel de la UE se limitó a autorizar a través de la Comisión ayudas al sector bancario. que necesitaban ser monitorizadas para impedir que afectaran negativamente a la competencia. Esto se debió en parte a la creencia de que se trataba de un asunto nacional y a la falta de reglas claras por parte de la UE.
} 

La crisis económica de 2008 demostró la necesidad de fortalecer la gobernanza económica y mejorar la coordinación de las políticas sociales entre los Estados miembros de la UE. En una unión de sistemas muy integrados, una mejor coordinación de las políticas económicas y sociales puede ayudar a prevenir discrepancias y contribuir a garantizar la convergencia y la estabilidad en el conjunto de la Unión y sus miembros. Se consideró necesaria una mayor coordinación de las políticas económicas en la UE a fin de abordar estos problemas e impulsar el crecimiento y la creación de empleo en el futuro. Los Estados miembros vieron entonces la necesidad de sincronizar el calendario de estos procedimientos para racionalizar el proceso y hacer más acordes los objetivos de las políticas nacionales presupuestarias, sociales, de crecimiento y de empleo, teniendo al mismo tiempo en cuenta los objetivos fijados a escala de la UE. Además, existía la necesidad de ampliar los ámbitos de supervisión y coordinación a políticas macroeconómicas y sociales más amplias.\footnote{Así, se revisó y reforzó el sistema de organismos y procedimientos de coordinación económica existentes en la UE en 2011, con la aprobación del \textit{Six-Pack}, en 2012 (con propuestas sobre la Unión bancaria y la creación del Mecanismo Europeo de Estabilidad), en 2013 (con la aprobación del \textit{Two-Pack}, paquete de dos medidas legislativas) y en 2014, con la creación del Mecanismo Único de Supervisión (MUS) y el Mecanismo Único de Resolución (MUR).
}

Por estos motivos, y en el marco de una reforma más amplia de la gobernanza económica y social de la UE, el Consejo Europeo decidió crear en 2010 el Semestre Europeo. El primer Semestre europeo tuvo lugar en 2011. Las bases jurídicas de este proceso son, en primer lugar, los artículos 121 y 148 del Tratado de Funcionamiento de la Unión Europea y, en segundo lugar, el denominado «paquete de seis medidas», seis actos legislativos que refonnaron el Pacto de Estabilidad y Crecimiento. Actualmente este proceso de coordinación se encuentra sujeto a revisión, con el objetivo de alcanzar un acuerdo entre todos los Estados miembros sobre un nuevo Semestre europeo reformado.

Cabe mencionar en este punto que la crisis económica desatada por la pandemia del COVID-$19$ supuso un reto sin precedentes para la Unión. Hasta 2019 se aplicó el Semestre europeo en su formato tradicional, pero las tensiones fueron cada vez más evidentes, hasta el punto de que la Comisión presentó en febrero de 2020 una propuesta para revisar el marco de Gobernanca económica de la Unión, abriendo el debate sobre su futuro. 

Sin embargo, el estallido de la crisis de COVID en marzo de 2020 dejó a un lado este debate para centrarse en la prioridad del momento: dar una respuesta coordinada y conjunta a nivel europeo. En este sentido, el formato del Semestre Europeo se mantuvo durante durante un periodo prolongado como una versión \textit{ad hoc}, adaptado a las circunstancias del momento y surgiendo de ésta la conocido activación de la claúsula general de escape.

Dos años más tarde, el debate se relanzó dos años después, con una propuesta presentada por la Comisión en noviembre de 2022 (Comunicación sobre la reforma de la Gobernanza económica de la Unión, con un foco especial sobre el proceso de supervisión fiscal europeo).\footnote{Los puntos más destacables de la propuesta fueron:
\begin{lnum}
    \item Mantener los límites del $-3\%$ del PIB para el déficit público y del $60\%$ del PIB para la deuda pública (ya que estos están establecidos en el TFUE, con poco margen de cambio) así como la cláusula de escape del PEC;
    \item Presentación por los Estados miembros de PFEMP. Estos Planes integrarán la doble dimensión de los dos procesos hasta ahora existentes: de desequilibrios macroeconómicos y de estabilidad fiscal, mediante sendas fiscales específicas por país de cuatro años, así como prioridades de reformas e inversiones públicas para garantizar una reducción gradual de la deuda y un crecimiento sostenible. Los Planes deberán ser valorados positivamente por la Comisión Europea y aprobados por elConsejo de la Unión Europea. Cabe la posibilidad de extender el periodo de aj uste a 7 años, si existe acuerdo sobre ello;
    \item El indicador operativo central será la senda multianual de gasto primario neto (\textit{interannual net primary expenditure path}), para situar la ratio deuda/PIB en una senda descendente y garantizar que el déficit no rebasa el límite del $-3\%$ del PIB a medio plazo;
    \item Se simplifican los indicadores cuantitativos para determinar el cumplimiento con la senda fiscal, que se centrará en el concepto de Techo de Gasto Neto Primario Nacionalmente Financiado (\textit{nationally financed net primary expenditure}). Esta definición implicará que las inversiones verdes y digitales financiadas con transferencias de fondos europeos no se tendrán en cuenta para el cumplimiento con el techo de gasto;
    \item Habrá mecanismos para sancionar el incumplimiento de los compromisos (una senda más dura de ajuste fiscal o incluso sanciones financieras aplicables sobre países de la zona euro). Se establecerían tres tipos de sanciones: (i) sanciones financieras, modestas pero frecuentes; (ii) sanciones de tipo reputacional, como intervenciones de los ministros en el Parlamento Europeo; o, (iii) cierta condicionalidad macroeconómica;
    \item Se dota de un rol más importante a las autoridades independiente fiscales nacionales (AIReF en España), en la evaluación de los programas y en su implementación. La Comisión y el Consejo tomarían su opinión en cuenta, pero no sería determinante.
\end{lnum}
}

Ese proceso de revisión ha culminado. Antes de profundizar en el funcionamiento del Semestre Europeo y, sobre todo, en el sistema de supervisión fiscal, conviene detenerse aquí a dar una panorámica de qué suponía el marco anterior, por qué se decidió reformarlo y en qué ha consistido, en esencia, el cambio.

Desde el Pacto de Estabilidad y Crecimiento de 1997 (reforzado por el \textit{Six-Pack}, 2011, el \textit{Fiscal Compact}, 2013, y el \textit{Two-Pack}, 2013) y la disciplina presupuestaria europea se articulaba sobre un conjunto de criterios numéricos concurrentes.\footnote{
Estas reglas eran el mantenimiento del déficit por debajo del $3\%$ del PIB y un nivel de deuda por debajo del $60\%$. Además, un ajuste anual mínimo del saldo estructural (0,5 p.p. del PIB) y, para los países con deuda superior al $60\%$, una regla de reducción de esa deuda al ritmo de 1/20 de la distancia respecto al umbral cada año. Merece la pena profundizar en el antiguo sistema para entender las razones que motivaron el cambio de paradigma. Su arquitectura de criterios múltiples nace de un compromiso político de mínimos entre Alemania (que exigía reglas numéricas duras y verificables) y el resto de socios, no de un diseño técnico óptimo; de ahí que la literatura la calificara reiteradamente de \textit{one-size-fits-all}, poco sensible a la heterogeneidad de situaciones fiscales dentro de la Unión.

Véase: \textit{Las nuevas reglas fiscales europeas: valoración e implicaciones para España}, Real Instituto Elcano (2024).
} 

Tres factores convergieron entre 2020 y 2023 que acabaron motivando su reforma. Primero, la cláusula de escape general activada en marzo de 2020 para dar respuesta a la pandemia (y prorrogada después por la crisis energética derivada de la guerra de Ucrania) debía desactivarse el 1 de enero de 2024, lo que habría supuesto la vuelta automática a unas reglas diseñadas para un contexto de deuda pública muy inferior al existente tras la pandemia\footnote{\textit{El Consejo acuerda reformar la gobernanza económica de la UE}, Fundación Galicia Europa (2024).}. Segundo, existía un consenso relativamente amplio en que el marco anterior era excesivamente complejo, procíclico en las fases de ajuste y poco compatible con la necesidad de proteger la inversión pública en las transiciones verde y digital, lección extraída precisamente de la insuficiente inversión tras la crisis de 2008. Tercero, la propia Comisión llevaba desde 2020 con las reglas fiscales en suspenso, lo que abría una ventana de oportunidad política para rediseñar el sistema en lugar de limitarse a reactivarlo.

En este contexto, la Comisión presentó sus orientaciones para la reforma en noviembre de 2022 y el paquete legislativo formal en abril de 2023. El acuerdo político se alcanzó en el Consejo ECOFIN el 20 de diciembre de 2023, durante la recta final de la Presidencia española del Consejo.\footnote{\textit{La Presidencia española en la UE: nueva gobernanza económica y financiera}. ICE, Revista de Economía, núm. 937 (2024)
}. Poco después, el trílogo con el Parlamento Europeo cerró el componente preventivo el 10 de febrero de 2024, el Consejo adoptó los tres textos legislativos definitivos el 29 de abril de 2024 y entraron en vigor al día siguiente, el 30 de abril de 2024.\footnote{Reglamento (UE) 2024/1263; Reglamento del Consejo de modificación del Reglamento (CE) n.º 1467/97; Directiva de modificación de la Directiva 2011/85/UE.

\textcolor{red}{Es interesante señalar en la exposición oral el papel de la Presidencia española (segundo semestre de 2023, con Nadia Calviño al frente de la negociación como vicepresidenta de la Comisión antes de su salto al BEI) como intermediaria del acuerdo político de diciembre, un dato que conecta este tema con el 3B36 y añade color institucional a una explicación que, de otro modo, es puramente normativa.}
}

\textcolor{red}{Un buen punto para lucirse ante el Tribunal es no presentar la reforma como un consenso pacífico, porque no lo fue. Se enfrentaron dos bloques con intereses contrapuestos: de un lado, Alemania y los países \textit{frugales} (Países Bajos, países nórdicos, Austria), que temían que la negociación bilateral entre Comisión y cada Estado miembro derivara en sendas de ajuste demasiado laxas, y que por ello lograron incorporar las salvaguardias numéricas mínimas (el ajuste de 0,5 p.p. si el déficit supera el 3\%, y los umbrales de reducción de deuda) como límite a esa discrecionalidad\footnote{Darvas, Z., Welslau, L. y Zettelmeyer, J. (2024), "The implications of the European Union's new fiscal rules", Policy Brief 10/2024, Bruegel.}; de otro, Francia, Italia y España, que defendían sendas de ajuste más graduales y vinculadas a reformas e inversiones. El resultado, señala Bruegel, exige en la práctica ajustes fiscales "ambiciosos" a los países muy endeudados, aunque menores que los que habría exigido el marco anterior. Y conviene citar también la crítica —más técnica, pero con enorme potencial de "lucimiento" doctrinal— de que el análisis de sostenibilidad de la deuda que sostiene todo el sistema es, en la práctica, un auténtica "caja negra": la propia réplica de Bruegel de esos cálculos exige más de tres mil líneas de código, lo que cuestiona si la reforma ha simplificado de verdad las reglas o simplemente ha trasladado la complejidad de la norma legal al modelo econométrico subyacente\footnote{CEPR/VoxEU, "Opening the black box of Europe's new fiscal rules", 2026.}. Añade además el Real Instituto Elcano una tercera crítica, distinta de las dos anteriores: el nuevo marco avanza poco en materia de aplicación efectiva y sanciones, por lo que su credibilidad última seguirá dependiendo, como con el marco anterior, de la voluntad política del Consejo\footnote{Real Instituto Elcano, "Las nuevas reglas fiscales europeas: valoración e implicaciones para España", 23/04/2024.}.}


\subsection{Marco actual: ¿Cómo se estructura la supervisión/coordinación de la Unión?}
Qué cambia, en síntesis (se detallará en II.III): un indicador operativo único, el gasto primario neto, sustituye a la pluralidad de criterios anterior; los Planes Fiscales-Estructurales a Medio Plazo, de cuatro a siete años, se convierten en la pieza central del sistema; y se introducen dos salvaguardias cuantitativas —de sostenibilidad de la deuda y de resiliencia del déficit— como contrapeso a la mayor discrecionalidad del nuevo modelo.

\subsubsection{Sistema de Supervisión Fiscal: Planes Fiscales-Estructurales a Medio Plazo (PFEMP)}
\textcolor{blue}{Este es el núcleo duro de la reforma, así que va a ser la sección más larga hasta ahora. Antes de escribirla: el documento original explica el **MTO** (objetivo presupuestario a medio plazo) como concepto central del brazo preventivo. Ese concepto **desaparece por completo** del nuevo marco — no se sustituye por otro MTO, se sustituye por una lógica distinta (senda de gasto derivada del análisis de sostenibilidad de la deuda). Es el cambio conceptual más importante de todo el tema, y es fácil pasarlo por alto si uno solo actualiza cifras en vez de repensar la arquitectura.}

El PEC se sigue articulando en dos vertientes, preventiva y correctiva, si bien el contenido de ambas ha sido reescrito por el Reglamento (UE) 2024/1263, el Reglamento del Consejo de modificación del Reglamento (CE) n.º 1467/97 y la Directiva de modificación de la Directiva 2011/85/UE, en vigor desde el 30 de abril de 2024.

**Una advertencia de método antes de seguir:** este epígrafe (II.III) es, con diferencia, el que más peso relativo va a tener en cualquier pregunta del Tribunal sobre este tema. Si el tiempo de repaso es limitado, es aquí donde hay que invertirlo primero, y no en el Semestre o el PDM.

\paragraph{Vertiente preventiva: Gasto Primario (neto)}
El sistema pivota íntegramente sobre un \textbf{indicador operativo único: el gasto primario neto}: gasto público excluido el pago de intereses, el componente cíclico del gasto en desempleo, el gasto cofinanciado con fondos europeos y las medidas puntuales, y neto de nuevas medidas de ingresos discrecionales\footnote{De esta forma, desaparece el objetivo presupuestario a medio plazo (MTO) específico de cada país, y con él desaparece la distinción entre saldo estructural observado y saldo estructural objetivo que dominaba el análisis desde 2011. La razón de fondo del cambio no es solo simplificar, sino cambiar el objeto de control de una variable no observable (el saldo estructural, que exige estimar el output gap y es, por definición, una variable estimada a partir de modelos macroeconómicos complejos basados en numerosas hipótesis y sujetos a gran incertidumbre y revisiones ex-post; crítica que también fue formulada respecto al Fiscal Compact. En definitiva, tomando como indicador de referencia el gasto primario neto se elige una \textbf{variable observable en tiempo real} (el gasto ejecutado). Se trata en esencia de un cambio de filosofía de control, no solo de fórmula.

Véase: \textit{El nuevo marco de gobernanza económica de la UE}, CaixaBank Research
}.

Además, la senda de gasto primario neto no la fija cada Estado miembro libremente: nace de un Análisis de Sostenibilidad de la Deuda (DSA) que realiza la Comisión, con proyección a diez años, combinando escenarios deterministas y estocásticos (tests de estrés), y que determina la trayectoria técnica de referencia para los Estados miembros con deuda superior al $60\%$ del PIB o déficit superior al $3\%$.\footnote{AIReF, "Observatorio de deuda", febrero 2024.} 

El Estado miembro incorpora esa trayectoria a su \textbf{Plan Fiscal-Estructural a Medio Plazo} (PFEMP), documento único que sustituye a los antiguos Programa de Estabilidad/Convergencia y Programa Nacional de Reformas, pudiendo negociar con la Comisión una extensión del periodo de ajuste de cuatro a sieta años a cambio de comprometer reformas e inversiones adicionales que mejoren la sostenibilidad de la deuda o el potencial de crecimiento.\footnote{Consilium, "Marco de gobernanza económica"; Portal MINECO, "Plan Fiscal y Estructural de medio plazo (2025-2028)"
}

Ahora bien, la trayectoria técnica debe respetar dos salvaguardias numéricas que actúan como suelo, no como objetivo:\footnote{\textit{Bruegel Policy Brief}, DARVAS, WELSLAU y ZETTELMEYER (10/2024).
}
\begin{la}
    \item \textbf{Sostenibilidad de la deuda}. La ratio de deuda debe reducirse, en promedio, al menos 1 punto porcentual de PIB al año si la deuda supera el $90\%$ del PIB, o 0,5 puntos si se sitúa entre el $60\%$ y el $90\%$; y,
    \item \textbf{Resiliencia del déficit}. Si el déficit estructural supera el $1,5\%$ del PIB (el margen de seguridad respecto al límite del $3\%$), el saldo estructural primario debe mejorar un mínimo de $0,4$ p.p. al año en un ajuste a cuatro años, o $0,25$ p.p. si el ajuste se extiende a siete.
\end{la}

Como se puede intuir, estas exigencia fueron introducidas a instancias de Alemania y los países frugales para evitar el riesgo de que esta que negociación bilateral Comisión-Estado derivara en sendas excesivamente laxas. Alemania y los nórdicos ganaron la batalla de las salvaguardias; Francia, Italia y España ganaron la batalla del horizonte temporal (4 a 7 años) y de la protección de la inversión.

El cumplimiento se controla mediante una \textbf{cuenta de control}, que registra anualmente las desviaciones (positivas y negativas) entre el gasto neto ejecutado y la senda comprometida, permitiendo cierta compensación intertemporal dentro de márgenes tasados. Este mecanismo es justamente el que ha permitido a España, en su Informe de Progreso 2025, incorporar como crédito a favor en la cuenta de control los cerca de $7.300\text{M}€$ ($0,5\%$ del PIB) de margen adicional derivado de la revisión de la senda por parte del Consejo, que podrá compensar desviaciones futuras. En esencia, un mecanismo de colchón que no existía en el sistema anterior, donde una desviación puntual podía activar directamente un Procedimiento de Desviación Significativa.\footnote{\textit{Informe de Progreso Anual 2025, Reino de España}. Ministerio de Hacienda.
}

Paralelamente, el marco refuerza el papel de las instituciones fiscales independientes nacionales (AIReF, en España), que evalúan no solo el cuadro macroeconómico subyacente (como ya ocurría) sino la propia senda de gasto y los Informes de Progreso Anual, con mayor independencia formal reconocida por la normativa europea.\footnote{BdE, Boletín Económico 2024/T4
}

\paragraph{Vertiente correctiva: Procedimiento de Déficit Excesico (PDE)}
El criterio del déficit, superior al $3\%$ del PIB, mantiene su activación del PDE sin cambios, así como la corrección mínima exigida de $0,5$ p.p. del PIB al año.

El criterio de la deuda sí se reescribe por completo. Desaparece la antigua regla de $1/20$ (reducción anual de un veinteavo de la distancia entre la deuda y el $60\%$ del PIB), que en la práctica resultaba inaplicable para los países más endeudados. ¿Por qué? Por ejemplo porque, aplicada sin cambios en 2024, la vieja regla habría exigido a Italia, con una deuda próxima al $135\%$ del PIB, una reducción de $3,75$ p.p. de PIB cada año.\footnote{Un ajuste de una magnitud sin precedente histórico en tiempos de paz, y que por sí solo explica por qué el consenso de reforma era prácticamente unánime aunque las razones concretas para reformarla difirieran entre países.

Véase: (i) \textit{Las nuevas reglas fiscales europeas: valoración e implicaciones para España}, Real Instituto Elcano (2024); y, (ii) \textit{Las nuevas reglas fiscales de la UE}, Agenda Pública (2024)
} En su lugar, el PDE por deuda se activa ahora de forma casuística: cuando la deuda supera el $60\%$ del PIB, el déficit no está próximo al equilibrio, y las desviaciones acumuladas en la cuenta de control superan determinados umbrales; la Comisión y el Consejo valoran, además, los retos de sostenibilidad de la deuda pública, el avance en reformas e inversiones y el aumento del gasto en defensa (novedad relevante y de plena actualidad).

Las sanciones se reducen y simplifican. La multa máxima pasa a ser del $0,05\%$ del PIB, acumulable cada seis meses hasta que el Estado miembro adopte medidas efectivas, frente al antiguo esquema de componente fijo ($0,2\%$) más componente variable (hasta $0,5\%$). Es aquí donde surgen mayores reticencias. Parte del mundo académico entiende que las nuevas reglas no avanzan mucho en el ámbito de la implementación práctica y de las sanciones, por lo que la credibilidad del sistema seguirá dependiendo, como en el marco anterior, de la voluntad política del Consejo más que del automatismo de la norma. Por tanto, permanece la cr´tica arrojada al PEC desde 2003, cuando Francia y Alemania inclumplieron sin consecuencias.

Por otro lado, se establece un régimen transitorio (2025-2027), a raíz del cual la Comisión puede, durante estos tres años, tener en cuenta el aumento de los pagos de intereses derivado de la subida de tipos al fijar la trayectoria correctora del PDE por deuda, para no penalizar dos veces a los países que ya sufren un mayor coste de financiación de su deuda.

Además, se mantienen la cláusula de escape general, para toda la Unión o la Eurozona en caso de recesión generalizada, y la clásusula nacional, para circunstancias excepcionales de un solo país, que ya existían en el marco anterior (COVID-$19$). No obstante, se precisa mejor su funcionamiento. En marzo de 2025, en el marco del plan ReArmar Europa/Preparación 2030, la Comisión activó una variante de la cláusula nacional específicamente pensada para el gasto en defensa, permitiendo a los Estados miembros superar temporalmente su límite de crecimiento de gasto neto para financiar rearme\footnote{\textit{El paquete de primavera del Semestre Europeo de 2025 ofrece orientaciones para impulsar la competitividad de la UE}, CDE Universidad de Granada
}. Prueba es de que, apenas un año después de su entrada en vigor, el marco de 2024 ya ha tenido que adaptarse a una circunstancia geopolítica que no estaba en el radar de sus diseñadores en 2023.

\subsubsection{Instrumentos de refuerzo: ¿qué más complementa?}
Junto al Pacto de Estabilidad y Crecimiento, la supervisión fiscal europea está caracterizada por otros acuerdos que refuerzan la actuación y ambición de dichas normas. 

\paragraph{Two-pack}
Uno de los más conocidos es el \textit{Two-Pack}, compuesto por dos reglamentos.

El primero de ellos, el Reglamento sobre disposiciones comunes para el seguimiento y la evaluación de los proyectos de planes presupuestarios y para la corrección del déficit excesivo de los Estados miembros de la zona del euro, establece calendarios presupuestarios comunes, así como normas para el seguimiento y evaluación de los Borradores de Planes Presupuestarios por parte de la Comisión. En caso de graves incumplimientos del PEC, la Comisión puede solicitar la revisión de los planes. Los Estados miembros sometidos a un PDE deben presentar un programa de asociación económica con las medidas y reformas necesarias para una corrección eficaz y duradera. Por otro lado tendríamos el Reglamento sobre el reforzamiento de la supervisión económica y presupuestaria de los Estados miembros de la zona del euro cuya estabilidad financiera experimenta o corre el riesgo de experimentar graves dificultades, que involucra al MEDE y establece normas para el reforzamiento de la supervisión, la ayuda financiera y la supervisión posterior al programa (hasta el reembolso de, al menos, el $75\%$ de la ayuda recibida)

Estos dos Reglamentos (472/2013 y 473/2013) no forman parte de los tres textos legislativos de la reforma de 2024 y permanecen en vigor sin modificación sustantiva.

\paragraph{Tratado de Estabilidad, Coordinación y Gobernanza en la UEM (2012): Pacto Presupuestario}
El pacto presupuestario (Fiscal Compact), incorporado ya al Derecho de la Unión.

En 2012 los Estados miembros firmaron el Tratado de Estabilidad, Coordinación y Gobernanza en la Unión Económica y Monetaria (TECG), cuyo Título III es el pacto presupuestario o Fiscal Compact. Este consagra la regla de oro del equilibrio presupuestario, con un límite mínimo del déficit estructural del $-0,5\%$ del PIB en la legislación nacional, preferiblemente de rango constitucional ($-1\%$ si la deuda es inferior al $60\%$ del PIB), y permite a los Estados miembros interponer recurso ante el TJUE contra otros Estados que no la hayan aplicado adecuadamente. La asistencia financiera del MEDE queda condicionada a la firma de este pacto por el Estado solicitante.\footnote{Todos los Estados miembros excepto Reino Unido y República Checa lo firmaron inicialmente; Croacia no lo hizo ni antes ni después de su adhesión en 2013; la República Checa lo ratificó en 2014.
}

El art. $16$ del TECG obligaba a incorporar su contenido al Derecho de la Unión en el plazo máximo de cinco años desde su entrada en vigor. Y, aquí sí, la reciente reforma ha jugado un papel fundamental: ese plazo se incumplió durante más de seis años y ha sido ésta la que ha ejecutado finalmente ese mandato, que incorporan expresamente al Derecho de la Unión la sustancia del Título III del TECG conforme a su art. $16$. No obstante, sobre esta incorporación existe controversia jurídica no cerrada ya que el propio Servicio Jurídico del Consejo, en su dictamen de septiembre de 2023, planteó tres preguntas sin respuesta unívoca: (i) si el paquete realmente incorpora la sustancia del pacto presupuestario en el sentido exigido por el art. $16$ TECG; (ii) qué consecuencias tiene esto sobre el régimen jurídico del propio TECG; y, (iii) si el art. $7$ del TECG (mayoría cualificada inversa para bloquear recomendaciones de la Comisión) sigue siendo aplicable una vez que el paquete de reforma lo sustituye en la práctica. Es decir, la regla de oro ya no depende solo de un tratado intergubernamental periférico al Derecho de la Unión, sino que forma parte del Derecho derivado de la UE, con las consecuencias que ello tiene en materia de control jurisdiccional por el TJUE y de primacía normativa.\footnote{La Directiva (UE) 2024/1265, que transpone este mandato junto con los requisitos generales sobre marcos presupuestarios nacionales, modifica sustancialmente el art. $7$ de la Directiva 2011/85/UE y suprime su art. $8$, alineando los requisitos de marco presupuestario a medio plazo nacional con el nuevo horizonte plurianual del PFEMP (cuatro a siete años, frente al mínimo de tres años anterior).
}

La reflexión crítica sobre la regla de oro que ya formulaba la doctrina respecto al Fiscal Compact sigue siendo plenamente aplicable, y gana relevancia ahora que esa regla pasa a formar parte del Derecho de la Unión propiamente dicho. En especial, el riesgo de restringir en exceso el déficit destinado a \textbf{financiar inversión} pública, con efectos adversos en equidad intergeneracional y eficiencia asignativa, la \textbf{dependencia de la sostenibilidad} de la deuda de variables distintas al mero equilibrio presupuestario (tipo de interés, crecimiento del PIB, nivel de deuda) y el \textbf{carácter no observable del saldo estructural}, sujeto a estimación y revisión ex post.\footnote{Crítica ya formulada en relación con el Fiscal Compact, aplicable con mayor razón tras su incorporación al Derecho derivado de la UE en 2024.
}

\subsubsection{Sistema de Supervisión Macroeconómica: Procedimiento de Desequilibrios Macroeconómicos (PDM)}
El Procedimiento de Desequilibrios Macroeconómicos (PDM) se creó como parte del \textit{Six-Pack}, vigente desde 2011, e integrado en el ciclo de supervisión multilateral del Semestre Europeo. Su finalidad es \textbf{vigilar y evitar la aparición de desequilibrios macroeconómicos potencialmente perjudiciales} que podrían afectar negativamente a la estabilidad económica de un Estado miembro concreto, de la zona del euro o de la Unión en su conjunto.\footnote{El PDM, a diferencia del PEC, no ha sido objeto de reforma legislativa en el paquete de 2024: los Reglamentos (UE) n.º 1176/2011 (vertiente preventiva) y n.º 1174/2011 (vertiente correctiva, sanciones) permanecen vigentes en su redacción original.

Banco de España, Boletín Económico 2024/T4, artículo 01: el nuevo marco de gobernanza mantiene al PDE y al PDM como los pilares básicos de su política de supervisión, si bien las nuevas normas aprobadas en abril suponen una revisión de facto del primero. 
}. 

Ante la ausencia de instrumentos compensatorios o mecanismos de ajuste de mercado, se hace necesaria la creación de herramientas de cooperación macroeconómica que complementen las existentes en materia de supervisión fiscal. El PDM nace con esta filosofía y tiene un doble objetivo: actuar como mecanismo preventivo de alerta temprana de desequilibrios, y realizar un análisis y seguimiento de los desequilibrios existentes con el objetivo de que se adopten medidas para su corrección.

El procedimiento se inicia en octubre-noviembre de cada año, con el inicio del Semestre Europeo, mediante la publicación del Informe sobre el Mecanismo de Alerta, cuyo eje central es un panel de catorce indicadores (MIP scoreboard), cada uno con un umbral de alerta, referidos tanto a desequilibrios externos (cuenta corriente, posición de inversión internacional neta, tipo de cambio efectivo real, cuotas de exportación, costes laborales unitarios) como internos (precios de la vivienda, flujo de crédito privado, deuda privada, deuda pública, tasa de desempleo, pasivos del sector financiero):

\textbf{INTRODUCIR}

Con base en este informe, la Comisión selecciona los Estados miembros a los que realizar un examen exhaustivo (In-Depth Review, IDR), documento analítico destinado a determinar y evaluar la gravedad de los desequilibrios detectados. Si se identifica un desequilibrio excesivo, la Comisión puede recomendar al Consejo la apertura de la vertiente correctora del PDM, el Procedimiento de Desequilibrio Excesivo (Excessive Imbalance Procedure, EIP), que puede acabar en sanciones de hasta el $0,1\%$ del PIB para los Estados miembros de la zona euro. El procedimiento concluye con la incorporación de sus conclusiones a las recomendaciones específicas por país del paquete de primavera. 

Recuadro 2. Situación de España en el Procedimiento de Desequilibrios Macroeconómicos (actualización)

España fue objeto de exámenes exhaustivos de forma ininterrumpida desde 2012, sin que se haya abierto nunca un EIP. El último IDR publicado en mayo de 2023 identificaba desequilibrios ligados a elevados niveles de deuda pública, privada y exterior, y a la situación del mercado laboral, con vulnerabilidades en remisión pero todavía presentes, un desempleo juvenil y de larga duración elevado, y una deuda pública en trayectoria descendente apoyada en el crecimiento del PIB nominal\footnote{Recuadro original del tema, IDR de mayo de 2023.}.

En el Informe sobre el Mecanismo de Alerta de noviembre de 2024, que inicia el ciclo del PDM para 2025, España queda por primera vez fuera del grupo de países seleccionados para examen exhaustivo, formado por Alemania, Chipre, Eslovaquia, Grecia, Hungría, Italia, Países Bajos, Rumanía, Suecia y, adicionalmente, Estonia por riesgo de nuevos desequilibrios\footnote{Servimedia, "Bruselas mantiene a España fuera de los diez países a examen exhaustivo en 2025 por 'desequilibrios'", 18/12/2024.}: la primera vez en trece ciclos consecutivos (2012-2024) que España no requiere examen exhaustivo, lo que la Comisión atribuye a la corrección sostenida de los desequilibrios exterior y de deuda privada iniciada tras 2012, unida al crecimiento económico sostenido desde 2022. 

No debe leerse, sin embargo, como ausencia de vulnerabilidad: la deuda pública española sigue entre las más elevadas de la zona euro en términos relativos, y el propio Consejo, en sus conclusiones sobre los exámenes exhaustivos de 2024 y 2025, ha insistido en que persisten riesgos importantes derivados de los elevados pasivos externos en varios Estados miembros y en que la vigilancia debe mantenerse.

La Comisión habría planteado en 2022 una revisión del PDM con un enfoque análogo al de las reglas fiscales, dando mayor peso a las tendencias de largo plazo, pero esa propuesta no llegó a materializarse en un texto legislativo. Lo que sí cambia es la interacción procedimental entre ambos procedimientos: cuando un Estado miembro presenta desequilibrios de naturaleza fiscal, la supervisión y la formulación de recomendaciones se canalizan a través del PEC mediante la apertura de un PDE, reservándose el PDM para los desequilibrios de naturaleza no estrictamente fiscal (competitividad, endeudamiento privado, sector inmobiliario, posición exterior).

\subsection{Etapas: ¿Cómo se organiza?}
El ciclo anual de coordinación de políticas comienza en noviembre, cuando la Comisión fija las prioridades para el año siguiente, y termina en octubre del año siguiente, cuando los gobiernos nacionales presentan los proyectos de planes presupuestarios teniendo en cuenta las recomendaciones de la UE adoptadas por el Consejo en verano. 

El Semestre Europeo sigue un ritmo anual. Los Estados miembros reciben asesoramiento a escala de la UE («orientaciones») y recomendaciones por separado (recomendaciones específicas por país) para sus políticas nacionales en materia presupuestaria y de reformas. Se invita a los Estados miembros a que tengan en cuenta esas recomendaciones al definir sus presupuestos para el ejercicio siguiente y al tomar decisiones estratégicas, en particular en los ámbitos de las políticas económicas, de empleo y educativas. Los Estados miembros también pueden recibir recomendaciones relativas a la corrección de desequilibrios macroeconómicos.

El Semestre Europeo abarca actualmente diferentes bloques de coordinación de las políticas económicas y sociales: las \textbf{reformas estructurales} (hoy articuladas fundamentalmente a través de los Planes de Recuperación y Resiliencia), las \textbf{políticas sociales y de empleo}, en consonancia con el pilar europeo de derechos sociales, las \textbf{políticas presupuestarias}, para garantizar la sostenibilidad de las finanzas públicas, y la \textbf{prevención de desequilibrios} macroeconómicos excesivos.

\subsubsection{Fase preparatoria: Paquete de Otoño (Nov.-Dic.)}
El ciclo del Semestre Europeo comienza cada año en noviembre, cuando la Comisión Europea publica el «paquete de otoño». Consta de varias partes:
un Estudio prospectivo anual sobre el crecimiento sostenible. El Estudio prospectivo anual sobre el crecimiento sostenible presenta la opinión de la Comisión sobre las prioridades económicas y sociales de la UE para el ejercicio siguiente. Se propone a los Estados miembros que las tengan en cuenta al elaborar sus políticas económicas para el ejercicio siguiente.
un Informe sobre el mecanismo de alerta. El Informe sobre el mecanismo de alerta examina la evolución macroeconómica en cada Estado miembro de la UE. El principal resultado del Informe sobre el mecanismo de alerta es una lista de los Estados miembros que podrían necesitar un examen exhaustivo. Se basa en un cuadro de indicadores que ayudan a controlar los desequilibrios externos, la competitividad y los desequilibrios internos en los Estados miembros. Estos indicadores se combinan con información adicional para evaluar el riesgo de desequilibrios. El hecho de que alguno de los indicadores sobrepase el umbral acordado en un país constituye la primera señal de posibles desequilibrios macroeconómicos. Esta información ayuda a la Comisión a decidir si lleva a cabo un examen exhaustivo. Basándose en el Informe sobre el mecanismo de alerta, la Comisión puede decidir llevar a cabo un examen exhaustivo de la situación en los Estados miembros en los casos en los que el riesgo de posibles desequilibrios macroeconómicos se considere elevado. Estos exámenes ayudan a determinar si existen posibles desequilibrios macroeconómicos y, en caso de que así sea, su naturaleza y alcance exactos. También permiten que la Comisión presente recomendaciones de actuación a los Estados miembros.
un proyecto de recomendación para la zona del euro. El proyecto de Recomendación del Consejo sobre la política económica de la zona del euro trata cuestiones clave para el funcionamiento de la zona del euro. El objetivo es integrar mejor la dimensión nacional y la de la zona del euro de la gobernanza económica de la UE.
un proyecto de informe conjunto sobre el empleo. El informe conjunto sobre el empleo ofrece una síntesis de los principales avances sociales y en materia de empleo en Europa, así como de las respuestas estratégicas de los Estados miembros. Supervisa los avances hacia los objetivos principales de la UE para 2030 en materia de empleo, capacidades y reducción de la pobreza, así como los avances hacia los objetivos nacionales para 2030 presentados por los Estados miembros. También analiza la aplicación de las orientaciones para las políticas de empleo de la UE. El informe conjunto sobre el empleo también supervisa los resultados de los Estados miembros en relación con el pilar europeo de derechos sociales, sobre la base del cuadro de indicadores sociales.


\paragraph{Octubre}
El 15 de octubre de cada año, los Estados miembros de la zona euro presentan sus Borradores de Planes Presupuestarios (DBP) del próximo año a la Comisión, que emite una Opinión en noviembre. 

Este trámite no ha cambiado con la reforma, que sigue siendo el único punto del calendario en el que la Comisión se pronuncia sobre el presupuesto anual propiamente dicho, y no sobre la trayectoria plurianual.

\paragraph{Noviembre: Paquete de Otoño}
En el marco del Paquete de otoño, la Comisión presenta el Dictamen sobre los proyectos de planes presupuestarios de los países de la zona euro, el Informe sobre el Mecanismo de Alerta (que pone en marcha el ciclo anual del PDM), el Informe conjunto sobre empleo, y las Recomendaciones específicas de política económica para la zona del euro en su conjunto.

Este bloque documental se mantiene prácticamente intacto.\footnote{\textit{La segunda parte del paquete de otoño del Semestre Europeo aborda los retos socioeconómicos de 2025}, Europe Direct Sevilla (2024)
}

\textcolor{red}{Vale la pena señalar un matiz analítico: el paquete de otoño ya no es solo diagnóstico, porque ahora también es el momento en el que (salvo excepción transitoria) se cierra el ciclo de los Planes Fiscales-Estructurales a Medio Plazo (PFEMP) de los Estados miembros que inauguran o renuevan un plan. En efecto, dado que el nuevo marco entró en vigor a finales de abril de 2024, el primer ciclo de PFEMP no pudo seguir el calendario ordinario: los 27 Estados miembros presentaron sus primeros planes en otoño de 2024 (España lo hizo en octubre), en lugar de en la primavera habitual.}

\subsubsection{Fase 1: Orientaciones políticas para la Unión}

Enero y febrero
El Consejo debate el Estudioprospectivo anual sobre el crecimiento sostenible y el Informe sobre el mecanismo de alerta y adopta unas Conclusiones.
También debate, modifica y aprueba la Recomendación sobre la política económica de la zona del euro, que posteriormente presenta al Consejo Europeo para que la debata en su reunión de marzo.
El Semestre Europeo tiene repercusiones para diversas políticas. El Consejo lo debate principalmente en sus formaciones Ecofin y EPSCO.
El Parlamento Europeo también debate el Estudio prospectivo anual sobre el crecimiento sostenible y puede publicar un informe por propia iniciativa. El Parlamento también participa en el Semestre Europeo a través del diálogo económico. Puede invitar al presidente del Consejo, a la Comisión y, cuando proceda, al presidente del Consejo Europeo o al presidente del Eurogrupo a tratar cuestiones relacionadas con el Semestre Europeo. Asimismo, se puede ofrecer a los distintos Estados miembros la oportunidad de participar en un intercambio de puntos de vista.

Marzo
El Consejo adopta el Informe conjunto sobre el empleo y adopta unas Conclusiones.
El Consejo Europeo debate la situación económica y del empleo en los Estados miembros y refrenda el proyecto de Recomendación del Consejo sobre la política económica de la zona del euro.

\subsubsection{Fase 2: Orientaciones específicas por país (Paquete de Primavera)}

\paragraph{Abril: Cambio más profundo}
Aquí es donde el ciclo se transforma por completo: \textcolor{red}{desaparecen, como documentos autónomos, los Programas de Estabilidad o Convergencia y los Programas Nacionales de Reforma que se presentaban cada mes de abril.} 

En su lugar, el nuevo marco establece que cada Estado miembro debe remitir a la Comisión, antes del 30 de abril, uno de estos dos documentos:
\begin{la}
    \item \textbf{Informe de Progreso Anual}. Si el PFEMP vigente sigue en curso —documento que da cuenta del cumplimiento de la senda de gasto neto acordada con el Consejo y del avance en los compromisos de reformas e inversiones\footnote{Ministerio de Hacienda, "Plan Fiscal y Estructural", hacienda.gob.es.}; o
    \item \textbf{PFEMP}. Si el periodo cubierto por el plan anterior ha finalizado.
\end{la}

España es un buen ejemplo de cómo funciona esto en la práctica: presentó su primer PFEMP (2025-2028) en octubre de 2024\footnote{Portal MINECO, Plan Fiscal y Estructural de medio plazo (2025-2028).}, y desde entonces ha remitido su Informe de Progreso Anual el 30 de abril de 2025 y el 30 de abril de 2026, en ambos casos confirmando el cumplimiento de la senda de gasto pactada\footnote{\textit{El Gobierno confirma el cumplimiento de las reglas fiscales europeas en 2025}, La Moncloa (2026)
}. \textcolor{red}{Un dato con mucho valor analítico para el "cante": el propio Informe de Progreso 2025 de España reconoce que, si se tienen en cuenta los márgenes de flexibilidad para el gasto en defensa activados por la Comisión (véase más abajo), el crecimiento del gasto neto español ($4,5\%$) se sitúa dentro de los límites del Plan, cuando sin esa flexibilidad estaría más cerca del límite\footnote{Portal MINECO, "2025-Informe-Anual-2026"; AIReF, "Informe de Seguimiento del PFEMP 2025-2028", mayo 2025.
}. Esto ilustra perfectamente cómo las "cláusulas de salvaguardia" no son una cláusula de escape general como la de la pandemia, sino un instrumento de ajuste fino, país por país y partida por partida.}
En consonancia con el nuevo marco de gobernanza económica de la UE, que se aplica desde 2024, cada Estado miembro presenta un plan fiscal-estructural nacional a medio plazo en el que establece su senda fiscal, así como sus prioridades de inversiones y reformas públicas que, en conjunto, garanticen una reducción sostenida y gradual de la deuda y un crecimiento sostenible e integrador, evitando una política fiscal procíclica.
Los planes tienen un período de aplicación de entre cuatro y cinco años. Permiten llevar a cabo un procedimiento de supervisión multilateral que también tenga en cuenta las tendencias a medio plazo. Esto ha mejorado el proceso de coordinación de las políticas, que anteriormente se examinaba cada año de forma aislada.
Los planes también establecen una senda de gasto neto: una trayectoria numérica plurianual para el gasto neto, que puede abarcar hasta siete años.
Sobre la base de una recomendación de la Comisión, el Consejo adopta una Recomendación que establece la senda de gasto neto. También puede aprobar las reformas e inversiones que sustenten cualquier ampliación de los períodos de ajuste.
Si un Estado miembro tiene un nuevo Gobierno, podrá presentar un plan fiscal-estructural nacional a medio plazo revisado.
En 2024, los Estados miembros deben preparar sus planes fiscales-estructurales en otoño. A partir de 2025, cada Estado miembro deberá preparar un informe anual de situación o un nuevo plan fiscal-estructural en abril (al final del período cubierto por el plan actual). El informe anual de situación describe la aplicación del plan fiscal-estructural nacional a medio plazo, que incluye la senda de gasto neto establecida por el Consejo, y las reformas e inversiones acometidas.

\paragraph{Mayo/junio: Paquete de primavera.} 
El antiguo mes de "mayo", dedicado a la evaluación de los PNR y los Programas de Estabilidad, se sustituye por el paquete de primavera, publicado normalmente entre mayo y junio. En él, la Comisión:

publica informes por país que hacen balance de la situación presupuestaria y evalúan los avances en la aplicación de las recomendaciones específicas de años anteriores;
compara el crecimiento previsto del gasto neto con el límite máximo fijado por el Consejo (o descrito en el propio plan nacional) — este es, en puridad, el nuevo criterio de cumplimiento, y sustituye por completo al antiguo análisis dual de saldo estructural y regla de gasto de referencia;
evalúa la posible existencia de desequilibrios macroeconómicos (PDM), pudiendo formular recomendaciones específicas al respecto; y
propone las recomendaciones específicas por país (REP/CSR), ya ancladas explícitamente en el cumplimiento de la senda de gasto y en la ejecución del PFEMP, y no en el logro de un objetivo de saldo estructural a medio plazo como antes.

\textcolor{red}{El primer paquete de primavera bajo el nuevo marco (2025) es un caso de estudio excelente porque muestra la reforma "en vivo" y con sus tensiones reales\footnote{Europe Direct Castilla y León, "Orientaciones del paquete de primavera del Semestre Europeo", junio 2025.}: (i) a Bélgica se le recomendó una nueva senda correctora tras detectarse que su crecimiento de gasto neto previsto superaba el límite fijado, aunque quedaba amparado por la cláusula nacional de salvaguardia; (ii) a Rumanía, en cambio, la Comisión concluyó que su desviación era tan significativa que recomendó al Consejo declarar que no había adoptado medidas efectivas para corregir su déficit excesivo —el primer test de credibilidad real del nuevo PDE por deuda—; y (iii) la Comisión activó, en el marco del plan "ReArmar Europa / Preparación 2030" (marzo de 2025), la cláusula nacional de salvaguardia de defensa, que permite a los Estados miembros superar temporalmente su tasa máxima de crecimiento de gasto neto para financiar gasto militar. Este último punto merece mención aparte en el epígrafe II.III cuando se explique el catálogo de salvaguardias, porque es la prueba de que el nuevo marco, pensado en 2023 sin Ucrania como variable central, ha tenido que incorporar sobre la marcha una excepción pensada para el rearme europeo.}

Como parte de su «paquete de primavera», la Comisión publica informes por país para todos los Estados miembros. Los informes hacen balance de la situación presupuestaria de los países y evalúan los avances realizados en la aplicación de las recomendaciones específicas por país de años anteriores.
El paquete de primavera también incluye exámenes exhaustivos de los desequilibrios macroeconómicos de los Estados miembros que puedan estar en riesgo de experimentar tales desequilibrios.
La Comisión también presenta recomendaciones específicas por país (REP). Estos documentos contienen orientaciones sobre políticas económicas, presupuestarias, de empleo y estructurales adaptadas para cada Estado miembro. Las REP recomiendan medidas que cada país debe adoptar para corregir los desequilibrios detectados. Estas recomendaciones pueden formularse al mismo tiempo que el examen exhaustivo o más adelante, junto con otras recomendaciones específicas por país.

\paragraph{Junio-Otoño}
\textcolor{yellow}{El Consejo aprueba las recomendaciones específicas para cada país, y las adopta oficialmente en julio, momento en el que se cierra el ciclo anual del Semestre Europeo.} 

Esto no cambia, aunque sí cambia su contenido: ahora las recomendaciones fiscales piden expresamente a los Estados miembros que respeten la tasa máxima de crecimiento del gasto neto fijada por el Consejo, y a los sujetos a un PDE que endurezcan su política fiscal para mantenerse dentro de la senda correctiva del procedimiento. En este momento pueden abrirse Procedimientos por Déficit Excesivo, Procedimientos de Desviación Significativa (respecto a la senda de gasto pactada, ya no respecto al saldo estructural) o el Procedimiento de Desequilibrios Macroeconómicos Excesivos, cuestiones que se analizarán con detalle más adelante.

Cabe mencionar en este punto que, a raíz de la crisis económica desatada por la pandemia del COVID- 19, se decidió en marzo de 2020 activar la conocida como cláusula general de escape, según la cual el Semestre europeo, y más concretamente el marco de supervisión fiscal europea, quedaba suspendido, con el objetivo de dotar de margen a los Estados miembros para adoptar las medidas de apoyo fiscal necesarias para hacer frente a una crisis sin precedentes. Dicha cláusula de escape se desactivó formalmente el 1 de enero de 2024, lo que —unido a la necesidad de dar continuidad a la flexibilidad ganada durante la pandemia— fue el detonante inmediato de la entrada en vigor del nuevo marco de gobernanza el 30 de abril de ese mismo año.

Tan pronto como la Comisión presente su paquete de primavera, comienza la supervisión presupuestaria multilateral. Este proceso dura varias semanas e implica un análisis y un debate en profundidad de todas las conclusiones pertinentes entre todos los Estados miembros y la Comisión. Como resultado de estos trabajos, el Consejo de la UE estará en condiciones de acordar las versiones finales de las recomendaciones específicas por país.
A continuación, el Consejo Europeo podrá debatir las recomendaciones finales, tras lo cual el Consejo podrá adoptarlas formalmente. Posteriormente se invita a los Estados miembros a aplicarlas.

\subsubsection{Fase 3: Implementación}
Los Estados miembros de la zona del euro tienen en cuenta las recomendaciones específicas por país a la hora de elaborar sus planes presupuestarios nacionales para el año siguiente. Deben presentar sus proyectos de planes presupuestarios al Consejo y a la Comisión a más tardar el 15 de octubre de cada año.
En el contexto del nuevo marco de gobernanza económica, en vigor desde el 30 de abril de 2024, las recomendaciones específicas por país tendrán un objetivo adicional. A partir de septiembre de 2024, todos los Estados miembros están obligados también a presentar planes fiscales-estructurales nacionales a medio plazo que abarquen un período de cuatro a cinco años, dependiendo de la duración habitual de la legislatura de cada Estado miembro. Este documento contiene los compromisos presupuestarios, de reformas y de inversiones del Estado miembro, así como una senda de gasto neto.
Algunos Estados miembros pueden intentar ampliar a siete años el periodo de ajuste de la senda de gasto neto, lo que requiere un esfuerzo presupuestario menos exigente por su parte. En tal caso, utilizarán las recomendaciones específicas por país que les correspondan como base para elaborar reformas e inversiones que puedan respaldar y justificar la ampliación.

\subsection{Valoración}
Ahora que conocemos la situación tras la reforma, ¿qué podemos decir al respecto?

Aunque el objetivo declarado de la reforma era \textit{simplificar} el sistema, un análisis reciente de CEPR/VoxEU sostiene que el análisis de sostenibilidad de la deuda que sustenta toda la arquitectura es, en realidad, una auténtica \textbf{caja negra}. De hecho, la propia réplica de estos cálculos por parte de Bruegel exige más de tres mil líneas de código\footnote{CEPR/VoxEU, "Opening the black box of Europe's new fiscal rules", 2026.}. 

La paradoja es que se ha sustituido una regla simple pero rígida (el saldo estructural) por un modelo más flexible pero opaco: la complejidad no ha desaparecido, se ha desplazado del texto legal al modelo econométrico. 


\begin{specialblock}{Situación de España}
    Recuadro 1. Situación de España en el PEC (actualización del recuadro original, que se mantiene como antecedente histórico hasta 2020): el déficit español descendió del 9,9\% del PIB en 2020 al 3,5\% en 2023, con objetivo del 3\% en 2024\footnote{Portal MINECO, "Plan Fiscal y Estructural de medio plazo 2025-2028".}; la deuda cayó desde su pico de 2021 hasta el 105\% del PIB en 2023, con objetivo del 102,5\% en 2024\footnote{Ídem.}; en 2025 el déficit se situó en el 2\% del PIB (excluyendo la DANA y otras partidas extraordinarias), sexto año consecutivo de cumplimiento del objetivo fijado por la Comisión\footnote{Portal MINECO, "2025-Informe-Anual-2026".}.

    **Recuadro 1 actualizado — España en el PEC.** El recuadro original (déficit del -10,1% en 2020, salida del PDE en 2019, etc.) sigue siendo válido como antecedente histórico, pero necesita una coda actualizada: el déficit español bajó del 9,9% (2020) al 3,5% en 2023, con previsión de alcanzar el 3% en 2024\footnote{Portal MINECO, "Plan Fiscal y Estructural de medio plazo 2025-2028".}; la deuda cayó del pico de 2021 hasta el 105% del PIB en 2023, con objetivo de 102,5% en 2024\footnote{Ídem.}; y en 2025 el déficit se situó en el 2% del PIB (excluyendo DANA y otras partidas extraordinarias), sexto año consecutivo de cumplimiento del objetivo fijado por la Comisión\footnote{Portal MINECO, "2025-Informe-Anual-2026".}.
\end{specialblock}






\clearpage
\section{Asisencia de la Unión: Instrumentos de asistencia financiera}
\puntosclave{
    Razón de ser e Historia: MEEF, FEEF y MEDE.
    
    Estructura actual y funciones, rol en la crisis del COVID (estigma político) 
    
    Ultima refonna del Tratado del MEDE pendiente de entrada en vigor a falta de la ratificación de Italia. 
    
    Breve mención al Mecanismo de ayuda a las balanzas de pagos.
}

Los dos apartados anteriores han descrito los mecanismos de coordinación y supervisión de las políticas económicas de los Estados miembros, cuya finalidad es fundamentalmente preventiva: detectar desequilibrios y corregirlos antes de que deriven en una crisis. Sin embargo, la experiencia de la crisis de deuda soberana de 2010-2012 demostró que la prevención no siempre es suficiente, y que la UEM necesitaba dotarse también de mecanismos institucionales de asistencia financiera para aquellos Estados miembros que, pese a la supervisión, llegaran a perder el acceso a los mercados de financiación en condiciones sostenibles.

Esta necesidad es particularmente acuciante en una unión monetaria como la Eurozona, donde los Estados miembros han cedido su política monetaria y, con ella, la posibilidad de actuar como prestamistas de última instancia de su propia deuda soberana en su propia moneda.\footnote{Este argumento es uno de los ejes centrales de la literatura AMO aplicada al diseño institucional del euro, y explica por qué la creación de un mecanismo de este tipo se consideró indispensable tras la crisis de deuda soberana
} Ante la ausencia de una capacidad fiscal centralizada de la Unión y de un prestamista de última instancia supranacional pleno, la respuesta institucional fue la creación, primero de soluciones de emergencia de carácter temporal y, en 2011, el Consejo Europeo decidió establecer un mecanismo permanente de resolución de crisis, el Mecanismo Europeo de Estabilidad (ESM por sus siglas en inglés), como instrumento de asistencia financiera a los Estados miembros de la zona euro con dificultades económicas graves. Se constituyó mediante un tratado intergubernamental, es decir, fuera del marco jurídico de la UE. 

De forma complementaria, para los Estados miembros de la UE que no pertenecen a la zona euro y que, por tanto, no pueden acceder al MEDE, existe desde 2002 el Mecanismo de Ayuda a las Balanzas de Pagos, con una lógica similar pero un ámbito subjetivo distinto.

Ambos instrumentos comparten una misma naturaleza: no son mecanismos de coordinación de políticas económicas, sino instrumentos de asistencia financiera condicionada, articulados sobre la lógica de la condicionalidad, que actúan como red de seguridad última cuando los mecanismos preventivos descritos en los apartados anteriores no han bastado para evitar una crisis de financiación.

\subsection{Mecanismo Europeo de Estabilidad (MEDE)}
La principal diferencia entre el MEDE y sus antecesores radica en su naturaleza jurídica y su alcance.\footnote{Por ejemplo, la FEEF era un organismo intergubernamental basado en acuerdos entre los países de la eurozona
} Tiene un carácter supranacional, cuyos accionistas son los Estados del Euro, otorgándole mayor capacidad de acción y una estructura más sólida y permanente.

La misión del MEDE es permitir a sus miembros evitar y superar las crisis financieras, asegurando la estabilidad financiera; misión que lleva a cabo mediante instrumentos financieros convencionales, con estricta concesionalidad (que se detalla en el MoU mediante reformas de ajuste), actuando como \textbf{prestamista de última instancia} cuando no pueden refinanciar su deuda a tipos sostenibles. \footnote{Tras cumplir con los requisitos de entrada en el euro, Croacia entra desde 2023 a formar parte del MEDE, como vigésimo miembro. La clave de contribución de cada estado miembro al MEDE está basada en la clave de capital del BCE.
} 

\subsubsection{Estructura de Gobernanza}
La estructura de gobernanza descansa en dos órganos principales: el \textbf{Consejo de Gobernadores}, en el que participan los ministros de economía de la zona euro y preside el presidente del Eurogrupo, y el \textbf{Consejo Directivo}, en el que participan los Directores (o Secretarios Generales) de los Tesoros. Asimismo, tiene también como observadores (sin voto) al Comisario de la cartera de \textit{Economía} (Asuntos Económicos) y al presidente del BCE. Paralelamente, el MEDE cuenta con un Consejo de Administración que asiste al Director Gerente en la actividad diaria del MEDE. 

Cuenta con un capital total suscrito de $700.000\text{M}€$, lo que le otorga una capacidad efectiva de préstamo de $500.000\text{M}€$, siendo acreedor preferente. Los préstamos se financian a través de empréstitos del MEDE en los mercados financieros y están garantizados por los accionistas (los Estados miembros de la zona euro).

Para otorgar mayor credibilidad a su actuación financiera, actúa en estrecha cooperación con otras instituciones internacionales, el FMI y el BCE, que también ofrecen apoyo en estas situaciones, de conformidad con sus objetivos y competencias respectivos.

Para recurrir al MEDE se debe haber aceptado el Pacto Presupuestario (\textit{Fiscal Compact}) mencionado en el apartado anterior. La asistencia financiera por parte del MEDE sólo se concede si se demuestra que es necesaria para salvaguardar la estabilidad financiera de la zona del euro en su conjunto y de los miembros del MEDE.

\subsubsection{Instrumentos del MEDE: Programas}
Para llevar a cabo su labor, el MEDE cuenta con un abánico variado de instrumentos financieros:
\begin{la}
    \item \textbf{Facilidades de crédito}. El primero, y más corriente, son los préstamos enmarcados en un programa de ajuste macroeconómico, cuyo objetivo es ayudar a los miembros del MEDE que han perdidó el acceso a los mercados, bien porque no pueden encontrar prestamistas, bien porque los costes de financiación repercutirían negativamente en la sostenibilidad de sus finanzas públicas;
    \item \textbf{Sindicación de deuda}. Los programas de compra de activos de deuda soberana en el mercado primario, permiten al MEDE participar en las operaciones de deuda sindicada por los miembros del MEDE, a precios de mercado para permitirles mantener o restablecer su relación con la comunidad inversora y reducir así el riesgo de una subasta fallida, aunque se limitará al $50\%$ del importe final emitido;
    \item \textbf{Programas de compra de títulos de deuda}. También pueden llevar a cabo operaciones similares en el mercado secundario, con el objetivo de apoyar el buen funcionamiento de los mercados cuando la prima de liquidez (en aumento) amenace la estabilidad financiera, ya sea con un programa de ajuste macroeconómico o sin él, y siempre y cuando la situación económica y financiera del país sea sólida;
    \item \textbf{Crédito precautorio}. Pueden abrir lineas de crédito precautoria, con el objetivo de apoyar la implementación de políticas y prevenir las situaciones de crisis. Se exige que las condiciones económicas sean sólidas y que permitan mantener un acceso continuo a la financiación del mercado, sirviendo únicamente para reforzar la credibilidad de sus proyecciones macroeconómicas;
    \item \textbf{Instrumentos de recapitalización: Sistema Bancario}. Como instrumento finalista también otorga préstamos para recapitalizar el sistema bancario, preservando la estabilidad del sector. El importe total disponible para este instrumento se limita a $60.000\text{M}€$ y se adhiere a la definición europea de los ámbitos materiales que comprende este sector (depósitos, sociedades financieras de cartera, sociedades financieras mixtas de cartera);
    \item \textbf{Instrumentos de recapitalización: Instituciones}. Por último, también podrían recapitalizar directamente las instituciones, contribuyendo a eliminar los riesgos de contagio originados en el sector financiero y dirigidos al soberano.
\end{la}

Los únicos tres programas que se han puesto en funcionamiento en toda la historia del MEDE son los programas de facilidades de crédito en el marco de programas de ajuste macroeconómico (Grecia y Chipre), y un préstamo para la recapitalización indirecta del Sistema Bancario concedido a España.

A raíz de la pandemia, el MEDE aprobó una nueva línea de crédito temporal: Apoyo a la Crisis Pandémica (Pandemic Crisis Support ECCL). A través de este instrumento se han financiado los costes relacionados con la sanidad, el tratamiento o la prevención de la pandemia, por un importe de hasta $240.000\text{M}€$ (hasta aproximadamente un $2\%$ del PIB de cada Estado miembro), y supone una mutación adaptativa de sus instrumentos ya que se omitió, a conciencia, cualquier otra condicionalidad que no fuese el destino material de los fondos obtenidos (gastos elegibles). De esta forma, los Estados miembros pueden financiar estos costes en condiciones homogéneas, repercutiéndose los costes con un margen de $10$ p.p. (más comisiones de servicio).\footnote{Aunque estuvo disponible para todos los sus miembros entre el 15 de mayo de 2020 y finales de 2022, y aunque no llevase aparejada ningún tipo de condicionalidad, ningún miembro hizo uso de los fondos. 

Aunque se han dado muchas explicaciones al respecto (estigma político, otros instrumentos como SURE o NGEU), la verdad es que no existe un consenso sobre las razones de su infrautilización. De todos modos, sí que exige una una revisión y reflexión profunda tanto sobre la situación, las funciones y su futuro.
}

\subsubsection{Reforma}
Tras un largo proceso de negociación entre los miembros del MEDE, se firmó en enero de 2021 la reforma del Tratado constitutivo, con el objetivo de reforzar el papel del MEDE como mecanismo de asistencia financiera, mejorando su capacidad de prevención y corrección de crisis en la eurozona. Sus principales elementos son: la provisión del backstop del MEDE al Fondo Único de Resolución mediante una línea de crédito renovable; el aumento de la efectividad de los instrumentos precautorios; la clarificación de las modalidades de cooperación entre el MEDE y la Comisión; y la mejora del marco de sostenibilidad de la deuda pública, simplificando el sistema de votación en las Cláusulas de Acción Colectiva y aumentando la transparencia de la metodología de análisis de sostenibilidad\footnote{ESM, "About ESM/ESM Reform Documents/ESM Treaty Amending Agreement".}.
 
 
Los principales puntos contemplados en la reforma son: 
• la provisión del ba ckstop del Mecanismo Europeo de Estabilidad al Fondo Único de Resolución a través de una línea de crédito renovable 
• aumento de la efectividad de los instrumentos precautorios, haciendo más transparentes y predecibles los criterios de elegibilidad 
• clarificación de las modalidades de cooperac ión entre el MEDE y la Comisión, tanto dentro como fuera de los programas de asistencia financiera 
• mejora del marco de sosten ibilidad de la deuda pública, simplificando el sistema de votación en las Cláusulas de Acción Colectiva, o aumentando la transparencia y previsibilidad de la metodología sobre el análisis de dicha sostenibilidad con la publicación de la metodología de las instituciones. 

La reforma requería ratificación unánime de los Estados miembros del MEDE. Todos los países la habían ratificado a finales de 2023 salvo Italia, único Estado miembro pendiente. El 21 de diciembre de 2023, la Cámara de Diputados italiana rechazó formalmente la ratificación, con 184 votos en contra, 72 a favor y 44 abstenciones.\footnote{Votaron en contra Hermanos de Italia (Meloni) y la Liga (Salvini); Forza Italia se abstuvo; el opositor Movimiento 5 Estrellas también votó en contra.
} Este resultado tiene un fuerte valor analítico y político: el rechazo italiano no fue un simple retraso técnico, sino una votación explícita en la que la coalición de gobierno se mostró dividida (Forza Italia, del PPE, se abstuvo, mientras Hermanos de Italia y la Liga votaron en contra por considerar la reforma un límite a la autonomía de gestión económica nacional), y en la que incluso el opositor M5S coincidió con el rechazo del Gobierno, acusando a Meloni de intentar aprobarla a escondidas en una votación anterior de 2021. El episodio refleja la disyuntiva, señalada por la literatura reciente, entre reformar el MEDE por la vía intergubernamental o transformarlo en un verdadero Fondo Monetario Europeo dentro del marco jurídico de la Unión, con gobernanza supranacional en lugar de intergubernamental.\footnote{Murua Astorquiza, O. (2025), "La controvertida reforma del mecanismo europeo de estabilidad (MEDE): un análisis jurídico constitucional", Revista Española de Derecho Europeo, núm. 93.
} El bloqueo italiano es, en este sentido, el ejemplo más claro y reciente de las limitaciones estructurales de la gobernanza puramente intergubernamental frente a la supranacional, tensión que recorre transversalmente todo el tema.

A fecha de esta actualización no consta una nueva votación de ratificación italiana, por lo que la reforma del Tratado del MEDE permanece sin entrar en vigor.

\subsubsection{Valoración}

El estigma político que ha caracterizado al MEDE desde la crisis financiera no se ha disipado: la reforma más orientada a completar la Unión Bancaria (el backstop al Fondo Único de Resolución) permanece bloqueada por motivos de soberanía fiscal percibida, más de una década después de la creación del mecanismo. 

La crisis de la Covid-19 volvió a mostrar, en cambio, que cuando la Unión Europea necesita movilizar una respuesta financiera de gran escala sin condicionalidad estigmatizante, tiende a crear instrumentos ad hoc (SURE, NGEU) antes que recurrir al MEDE, cuestión que se desarrolla en la siguiente sección.

\subsection{Mecanismo de Ayuda a la Balanza de Pagos}
Desde febrero de 2002, la UE ofrece ayuda a los países de la UE no pertenecientes a la zona del euro. La ayuda a la balanza de pagos adopta la forma de préstamos a medio plazo condicionados a la aplicación de políticas encaminadas a resolver los problemas económicos subyacentes. Normalmente, la ayuda de la UE a la balanza de pagos se ofrece en cooperación con el Fondo Monetario Internacional (FMI) y otras instituciones internacionales o países. Se ha prestado este tipo de asistencia financiera a Hungría, Letonia y Rumanía.

El país que solicita la ayuda y la Comisión firman un Memorándum de Entendimiento (MoU) en el que se exponen las condiciones de política económica que deben cumplirse antes de que se liberen los fondos, y un acuerdo de préstamo, en el que se detallan el proceso de endeudamiento y las condiciones financieras del préstamo. Este MoU contiene condiciones de política económica que suelen implicar medidas para garantizar la solidez de las finanzas públicas y la estabilidad del sector financiero, reformas estructurales para mejorar la competitividad y el crecimiento económicos, y salvaguardias contra el fraude. La ayuda financiera se desembolsa por etapas unavez verificado el cumplim iento de las condiciones del MoU.



\clearpage
\section{COVID-$19$: Política Económica de la Unión}
\puntosclave{
    Medidas inmediatas a nivel europeo (activación de la cláusula general de escape, nuevo marco de ayudas de Estado, tlexibilización en el uso de fondos europeos) 
    
    Medidas de medio plazo (SURE, nueva línea de crédito del MEDE y Fondo Paneuropeo de Garantías del BEi) 
    
    Medida sin precedentes: creación de Next Generation EU y del Mecanismo de Recuperación y Resiliencia 
    
    Valoración de la respuesta europea y debate sobre la creación de una capacidad central de estabilización
}

La crisis económica desatada por la pandemia del COVID-19 supuso un reto sin precedentes para la Unión Europea. Sin embargo, y a diferencia de lo ocurrido en la crisis financiera de 2008, la Unión Europea se movilizó para dar una respuesta conjunta que logró materializar en unos pocos meses. Esto resultó esencial para garantizar la igualdad de oportunidades de los distintos países en la fase de recuperación, independientemente de su capacidad individual de apoyar sus economías, evitando así la creación de nuevas divergencias.

Esta respuesta coordinada a nivel europeo ha supuesto un gran avance en el proceso de integración europeo: por primera vez en la historia, la Comisión europea se endeudó en los mercados en semejante magnitud, 750.000 millones de euros, 52% de los cuales, destinados a financiar transferencias directas a los Estados miembros para hacer frente, no solo al reto de la pandemia, si no a la transformación verde y digital de sus economías.

Si la contundente respuesta fiscal en un primer momento tuvo un carácter fundamentalmente nacional, desde el principio quedó en evidencia que la respuesta debía ser coordinada, sobre todo a nivel europeo. Frente a un shock exógeno común, es fundamental preservar la igualdad de oportunidades en todos los países europeos para evitar fragmentar el mercado interior. Además, las importantes interconexiones entre las economías europeas implican que, si una región se quedara atrás en la fase de recuperación, las consecuencias serían negativas para el conjunto de la UE.

Lo cierto es que, aunque el shock del Covid fue exógeno, el impacto fue asimétrico. No sólo en términos de la gravedad de la pandemia, sino también en la fase de recuperación, dependiendo de la especialización sectorial de las economías. Esto requirió que la respuesta coordinada se centrara por tanto en las regiones y sectores más afectados, contribuyendo efectivamente a paliar el impacto de la crisis. La acción coordinada a nivel europeo estaba justificada, en particular para restablecer la igualdad de condiciones y evitar acrecentar las divergencias entre los Estados miembros de la UE.

La respuesta europea a la crisis, podría resumirse en los siguientes "paquetes": (i) la respuesta inmediata de las autoridades europeas para apoyar las medidas aplicadas a nivel nacional (activación de la cláusula general de escape del Pacto de Estabilidad y Crecimiento; nuevo marcoeuropeo temporal de ayudas de Estado; tlexibilización en el uso de los fondos europeos27, etc)28 ; (ii) las medidas a corto plazo acordadas en el seno del Eurogrupo, que se componen de tres instrumentos: la línea de crédito del Mecanismo Europeo de Estabilidad, el instrumento SURE29 y el Fondo Paneuropeo de Garantías del Banco Europeo de Inversiones; y, (iii) el paquete de recuperación europeo acordado por los líderes políticos en el Consejo Europeo del 21 de julio, un paquete integral de 1.824,3 mil millones de euros que combina el Marco Financiero Plurianual 2021 -2027 y fondos adicionales obtenidos mediante emisión de deuda a través de un nuevo instrumento temporal para la recuperación, Next Generation EU.

Para atender a esa necesidad de proporcionar una re spuesta europea adecuada, el 9 de abril de 2020 el Eurogrupo alcanzó un acuerdo sobre tres mecanismos para responder a las necesidades a corto plazo derivadas de la crisis, basados en préstamos y garantías, además de apuntar en ese acuerdo a la puesta en marcha de un fondo de recuperación que se analizará más adelante. Los tres instrumentos suponen un paquete de 540.000 millones de euros que crean una triple red de seguridad para trabajadores (instrumento SURE de la Comisión), empresas (Fondo Paneuropeo de Garantías del Banco Europeo de Inversiones) y gobiernos (línea de crédito del Mecanismo Europeo de Estabilidad, explicada anteriormente).

Cabe mencionar las medidas adoptadas con prontitud por las autoridades monetarias, medidas destinadas a contener los costes de fi nanciación y preservar la liquidez. En marzo de 2020, el Banco Central Europeo amplió los programas de compras de activos existentes y adoptó, de forma temporal, el Programa de Compras de Emergencia para la Pandemia30 (PEPP, por sus siglas en inglés), dotado con 1,85 billones de euros31• Igualmente, para apoyar la liquidez en la economía real, se flexibilizaron ciertos parámetros en los programas de concesión de liquidez a las entidades de crédito, desarrollando además nuevas .operaciones de financiación a · plazo más largo (PELTRO).

Garantías del BEI

A través del Fondo Paneuropeo de Garantías gestionado por el BEi, los Estados miembros aportan garantías por un importe de 25.000 millones de euros con el objetivo de movilizar financiación por un importe de hasta 200.000 millones de euros. Las empresas de toda la UE pueden así acceder a una garantía de la misma calidad crediticia (AAA), permitiendo paliar en parte el posible problema de fragmentación del mercado interior por la distinta capacidad de las administraciones públicas de proporcionar apoyo a sus empresas.

SURE

A través del mecanismo SURE, la Comisión otorga préstamos a los Estados miembros para financiar esquemas de mantenimiento del empleo o medidas simi lares, por un importe máximode 100.000 millones de euros (con un límite de 60.000 millones de euros para los tres mayoresreceptores). Para ello, la Comisión se endeuda en el mercado y traslada las condiciones definanciación a los Estados miembros, lo que permite en definitiva homogeneizar las condiciones en las que los países financian estas medidas, solventando en parte el problema de fragmentación financiera. Según datos del último desembolso (marzo 2022) disponible hasta la fe cha33, el instrumento SURE ha proporcionado asistencia financiera por un montante total de 91.800 millones de euros, siendo los principales beneficiarios del instrumento SURE Italia (ha recibido 27.440 millones de euros) y España (ha recibido $21.324\text{M}€$).

NGEU

Si bien la triple red de seguridad fue un gran avance, el verdadero gran hito de esta respuesta europea coordinada se materializó en la creación de Ne.xi Generation EU, un nuevo instrumento temporal basado en el artículo 122 del Tratado de Funcionamiento de la UE, que permite a la Comisión endeudarse en los mercados por un valor de 750.000 millones de euros, emitiendo bonos con vencim ientos comprendidos entre 2028 y 2058. Se trata de un instrumento de emergencia único, puesto en marcha durante un período temporal y utilizado exclusivamente para la respuesta a la crisis y las medidas de recuperación. Se acordó que los Estados miembros no realizarían contribuciones adicionales al presupuesto de la UE durante el período 202 1-2027, ya que los primeros reembolsos tendrían lugar a partir del próximo marco financiero plurianual y durante un período prolongado (de 2028 a 2058). Además, la Comisión puso de manifiesto su intención de proponer nuevos recursos propios para hacer frente al repago de esta deuda, minimizando el impacto en las contribuciones nacionales. En cuanto al uso de los recursos de • NGEU, éstos se canalizan a través del Marco Financiero Plurianual 2021-2027, 390.000 millones de euros en forma de subvenciones a los Estados miembros y 360.000 millones de euros en préstamos, a través de su instrumento fundamental, el Mecanismo de Recuperación y


Resil iencia34.

El Mecanismo de Recuperación y Resiliencia proporc iona apoyo financiero a gran escala a las inversiones y reformas previstas en los planes nacionales de recuperación y resiliencia presentados por los Estados Miembros (como el Plan de Recuperación, Transformación y Resiliencia español). El instrumento, que se basa en el artículo 175 del TFUE (base legal de la política de cohesión), apoya la recuperación de las regiones más afectadas por la crisis y se centra así mismo especialmente en la transición verde y digital y en las prioridades identificadas en el Semestre Europeo (de hecho para el Semestre europeo 2023 se ha decidido ligar de fonna más estrecha el proceso del Mecanismo de Recuperación y Resiliencia al proceso del Semestre  europeo, para garantizar una mayor coherencia entre ambos35).

El acuerdo alcanzado por el Consejo europeo por el que se crea Ne.xt Generation E U ha supuesto un paso importante, ya que, por primera vez, la Comisión emite deuda en mercados en semejantes cantidades para financiar programas del presupuesto comunitario, apoyando a aquellos Estados miembros más afectados por la crisis. No obstante, Ne.xt Generation EU es un instrumento temporal, sin vocación de permanencia dada su base legal, el artículo 122 del Tratado de Funcionamiento de la Unión Europea, por lo que la idea de crear una verdadera capacidad central de estabilización quedaría descartada por el momento.


\subsubsection{Valoración}











\clearpage
\section{Conclusión} 


\subsection{Recapitulación}


\subsection{Extensiones y relación con otras partes del Temario}



\subsection{Opinión}

\subsubsection{Idea final (Salida o cierra)}

\clearpage
\pagenumbering{gobble}
\MyBibliography



\MyBibItem{DAWID y GATTI}{Chapter 2 - Agent-Based Macroeconomics}{2018}{https://doi.org/10.1016/bs.hescom.2018.02.006}


% ANEXOS
\clearpage
\anexo{\textbf{Preguntas Test}}
%\input{\TemaCode/questions.tex} 



\end{document}